% !TEX root = main.tex
\section{From raw text to a trained language model}

\subsection*{Notation}

\begin{tabular}{c l}
  \(B\) & batch size \\
  \(T\) & sequence length \\
  \(C\) & embedding/model dimension \\
  \(H\) & number of attention heads \\
  \(D\) & head size, \(D=C/H\) \\
  \(V\) & vocabulary size
\end{tabular}

\subsection*{Training and generation pipeline}

\begingroup
\setlength{\parskip}{0pt}
\small\ttfamily
\begin{longtable}{@{}p{0.67\textwidth}@{}l@{}}
RAW TEXT & \\
\hspace*{2em}\(\downarrow\) & \\
Build vocabulary & \([V]\) \\
\hspace*{2em}\(\downarrow\) & \\
Encode characters/tokens into IDs & \([T]\) \\
\hspace*{2em}\(\downarrow\) & \\
Create training batches & \([B,T]\) \\
\hspace*{2em}\(\downarrow\) & \\
Token embeddings & \([B,T,C]\) \\
\hspace*{2em}+ & \\
Position embeddings & \([T,C]\) \\
\hspace*{2em}\(\downarrow\) & \\
Representations & \([B,T,C]\) \\
\hspace*{2em}\(\downarrow\) & \\
Transformer Block & \([B,T,C]\) \\
\hspace*{2em}|-- LayerNorm & \([B,T,C]\) \\
\hspace*{2em}|-- Multi-Head Self-Attention & \\
\hspace*{2em}|\hspace*{3em}|-- Query & \([B,H,T,D]\) \\
\hspace*{2em}|\hspace*{3em}|-- Key & \([B,H,T,D]\) \\
\hspace*{2em}|\hspace*{3em}|-- Value & \([B,H,T,D]\) \\
\hspace*{2em}|\hspace*{3em}|-- Attention scores & \([B,H,T,T]\) \\
\hspace*{2em}|\hspace*{3em}|-- Scaling & \([B,H,T,T]\) \\
\hspace*{2em}|\hspace*{3em}|-- Causal masking & \([B,H,T,T]\) \\
\hspace*{2em}|\hspace*{3em}|-- Softmax / attention weights & \([B,H,T,T]\) \\
\hspace*{2em}|\hspace*{3em}|-- Dropout & \([B,H,T,T]\) \\
\hspace*{2em}|\hspace*{3em}|-- Weighted values & \([B,H,T,D]\) \\
\hspace*{2em}|\hspace*{3em}`-- Combine heads & \([B,T,C]\) \\
\hspace*{2em}|-- Residual connection & \([B,T,C]\) \\
\hspace*{2em}|-- LayerNorm & \([B,T,C]\) \\
\hspace*{2em}|-- Feed-Forward Network & \\
\hspace*{2em}|\hspace*{3em}|-- Expand & \([B,T,4C]\) \\
\hspace*{2em}|\hspace*{3em}`-- Project back & \([B,T,C]\) \\
\hspace*{2em}`-- Residual connection & \([B,T,C]\) \\
\hspace*{2em}\(\downarrow\) & \\
Repeat Transformer blocks & \([B,T,C]\) \\
\hspace*{2em}\(\downarrow\) & \\
Final LayerNorm & \([B,T,C]\) \\
\hspace*{2em}\(\downarrow\) & \\
Language-model head & \([B,T,V]\) \\
\hspace*{2em}\(\downarrow\) & \\
Logits & \([B,T,V]\) \\
\hspace*{2em}\(\downarrow\) & \\
Flatten for loss & \\
\hspace*{2em}|-- Logits & \([B\mathbin{\times}T,V]\) \\
\hspace*{2em}`-- Targets & \([B\mathbin{\times}T]\) \\
\hspace*{2em}\(\downarrow\) & \\
Cross-entropy loss & \([\text{scalar}]\) \\
\hspace*{2em}\(\downarrow\) & \\
Backpropagation & \\
\hspace*{2em}\(\downarrow\) & \\
AdamW updates parameters & \\
\hspace*{2em}\(\downarrow\) & \\
Repeat training & \\
\hspace*{2em}\(\downarrow\) & \\
Generation & \\
\hspace*{2em}|-- Context tokens & \([B,T]\) \\
\hspace*{2em}|-- Model logits & \([B,T,V]\) \\
\hspace*{2em}|-- Latest logits & \([B,V]\) \\
\hspace*{2em}|-- Softmax probabilities & \([B,V]\) \\
\hspace*{2em}|-- Sample next token & \([B,1]\) \\
\hspace*{2em}`-- Append token & \([B,T+1]\) \\
\hspace*{2em}\(\downarrow\) & \\
Repeat &
\end{longtable}
\endgroup

\clearpage
\section{Formula, purpose, and tiny examples}

\subsection{Build vocabulary \( [V] \)}

\textbf{What it does:} Finds every unique token the model can understand and
predict.

\[
  V = \lvert\text{unique tokens}\rvert
\]

\textbf{Example:}
\begin{verbatim}
text = "hello"
tokens = {h,e,l,o}
V = 4
\end{verbatim}

\subsection{Encode text into token IDs}

\textbf{What it does:} Converts text into integers.

\[
  \text{token} \longrightarrow \text{token ID}
\]

\textbf{Example:}
\begin{verbatim}
h -> 0
e -> 1
l -> 2
o -> 3

"hello" -> [0,1,2,2,3]
\end{verbatim}

\subsection{Create training batches \( [B,T] \)}

\textbf{What it does:} Creates several sequences for training at once.

\[
  X \in \mathbb{Z}^{B\times T},
  \qquad
  Y_{b,t} = X_{b,t+1}
\]

The targets are shifted by one token. \\

\break

\textbf{Example:}
\begin{verbatim}
Input:   h e l l
Target:  e l l o

32 = B = batch size
8  = T = sequence length

inputs.shape  = [32,8]
targets.shape = [32,8]
\end{verbatim}

Conceptually, the input batch could look like this:
\begin{verbatim}
inputs = [
  [h, e, l, l, o, <space>, w, o],  <- sequence 1
  [e, l, l, o, <space>, w, o, r],  <- sequence 2
  [l, l, o, <space>, w, o, r, l],  <- sequence 3
  ...
  32 sequences total
]
\end{verbatim}

Each row contains eight tokens. Therefore, the batch shape is
\[
  \underbrace{32}_{\substack{\text{sequences}\\\text{in the batch}}}
  \;\times\;
  \underbrace{8}_{\substack{\text{tokens in}\\\text{each sequence}}}
  \;=\;[32,8].
\]

\subsection{Token embeddings \( [B,T,C] \)}

\textbf{What it does:} Converts each token ID into a learned vector.

\[
  E_{\text{token}}(x_t) \in \mathbb{R}^{C},
  \qquad
  [B,T] \longrightarrow [B,T,C]
\]

\textbf{Example:} Token \texttt{h}, with ID 3, might become
\[
  [0.2,-0.7,0.4,\ldots].
\]
If \(C=32\), each token gets 32 numbers.

\begin{quote}
  ``What token is this?''
\end{quote}

\subsection{Position embeddings \( [T,C] \)}

\textbf{What it does:} Gives the model information about where each token
appears.

\[
  E_{\text{position}}(t) \in \mathbb{R}^{C}
\]

The initial representation is
\[
  x_t = E_{\text{token}}(\text{token}_t)
        + E_{\text{position}}(t).
\]

\textbf{Example:}
\begin{verbatim}
first "l":
token_embedding(l) + position_embedding(2)

second "l":
token_embedding(l) + position_embedding(3)
\end{verbatim}

The same token at a different position has a different representation.

\begin{quote}
  ``Where is this token?''
\end{quote}

\subsection{Initial representation \( [B,T,C] \)}

After adding both embeddings,
\[
  X = E_{\text{token}} + E_{\text{position}}.
\]

\textbf{Example:} Shape \([32,8,32]\) means 32 sequences, 8 tokens per
sequence, and 32 features per token.

\subsection{Layer Normalization}

\textbf{What it does:} Keeps each token's feature values numerically stable.
For one token vector,
\[
  \mu = \frac{1}{C}\sum_{i=1}^{C}x_i,
  \qquad
  \sigma^2 = \frac{1}{C}\sum_{i=1}^{C}(x_i-\mu)^2,
\]
\[
  \widehat{x}_i = \frac{x_i-\mu}{\sqrt{\sigma^2+\epsilon}},
  \qquad
  y_i = \gamma_i\widehat{x}_i+\beta_i.
\]

\textbf{Example:}
\begin{verbatim}
before: [2,4,6,8]
after:  approximately [-1.34,-0.45,0.45,1.34]
\end{verbatim}

\begin{quote}
  ``Keep the numerical scale stable.''
\end{quote}

\subsection{Query, Key, and Value}

For an input \(X\in\mathbb{R}^{B\times T\times C}\), the model creates
\[
  Q=XW_Q, \qquad K=XW_K, \qquad V=XW_V.
\]

Initially, \(Q\), \(K\), and \(V\) each have shape \([B,T,C]\).

\begin{center}
\begin{tabular}{@{}ll@{}}
Query & ``What am I looking for?'' \\
Key   & ``What information do I contain?'' \\
Value & ``What information should I send?''
\end{tabular}
\end{center}

Suppose the input is:
\begin{verbatim}
the cat sat
\end{verbatim}

Its tensor shape is
\[
  X:[B,T,C]=[1,3,4].
\]
This means there are:
\begin{itemize}
  \setlength{\itemsep}{0pt}
  \setlength{\parsep}{0pt}
  \item 1 sequence,
  \item 3 tokens, and
  \item each token represented by 4 numbers
\end{itemize}

For example,
\begin{align*}
  x_{\text{the}} &= [1.0,0.2,0.1,0.7],\\
  x_{\text{cat}} &= [0.3,1.1,0.8,0.2],\\
  x_{\text{sat}} &= [0.5,0.4,1.2,0.9].
\end{align*}
Thus, omitting the outer batch dimension for readability,
\[
  X=
  \begin{bmatrix}
    1.0 & 0.2 & 0.1 & 0.7\\
    0.3 & 1.1 & 0.8 & 0.2\\
    0.5 & 0.4 & 1.2 & 0.9
  \end{bmatrix}
  \begin{matrix}
    \leftarrow \text{the}\\
    \leftarrow \text{cat}\\
    \leftarrow \text{sat}
  \end{matrix}.
\]

The model starts with the same representation \(X\), but learns three
different projection matrices, \(W_Q\), \(W_K\), and \(W_V\), and computes
\[
  Q=XW_Q, \qquad K=XW_K, \qquad V=XW_V.
\]
The original token vectors are therefore transformed in three different ways. \\

With the linear layers
\begin{lstlisting}[language=Python]
self.query = nn.Linear(C, C, bias=False)
self.key   = nn.Linear(C, C, bias=False)
self.value = nn.Linear(C, C, bias=False)
\end{lstlisting}
the shapes remain
\[
  Q:[B,T,C], \qquad K:[B,T,C], \qquad V:[B,T,C].
\]

For example, the token \texttt{sat} might begin as
\[
  x_{\text{sat}}=[0.5,0.4,1.2,0.9]
\]
and become
\begin{align*}
  q_{\text{sat}} &= [0.8,0.1,0.7,0.2],\\
  k_{\text{sat}} &= [0.2,0.9,0.4,0.6],\\
  v_{\text{sat}} &= [1.1,0.3,0.8,0.5].
\end{align*}

\subsubsection*{Query: ``What am I looking for?''}

When processing \texttt{sat}, its query might represent the question:
\begin{quote}
  \emph{``I am a verb. Which previous token is likely the subject performing
  this action?''}
\end{quote}
The model uses this query to search the other tokens for relevant information.

\subsubsection*{Key: ``What information do I contain?''}

Each token has a key that advertises the kind of information it contains. For
example,
\begin{center}
\begin{tabular}{@{}ll@{}}
\(k_{\text{the}}\) & ``I am mostly a determiner.'' \\
\(k_{\text{cat}}\) & ``I am a noun and a possible subject.'' \\
\(k_{\text{sat}}\) & ``I am a verb.''
\end{tabular}
\end{center}

The query for \texttt{sat} is compared with every key. Example dot products
might be
\begin{align*}
  q_{\text{sat}}\mathbin{\cdot}k_{\text{the}} &= 0.3,\\
  q_{\text{sat}}\mathbin{\cdot}k_{\text{cat}} &= 2.4,\\
  q_{\text{sat}}\mathbin{\cdot}k_{\text{sat}} &= 0.4.
\end{align*}
The strongest match is \texttt{cat}, so attention treats \texttt{cat} as
especially relevant when updating the representation of \texttt{sat}.

\subsubsection*{Value: ``What information should I send?''}

The key helps determine whether a token is relevant. Once the attention
weights have been calculated, the value supplies the actual information to
transfer. Suppose the attention weights for \texttt{sat} are
\[
  \text{the}:10\%, \qquad
  \text{cat}:70\%, \qquad
  \text{sat}:20\%.
\]
Then the output for \texttt{sat} is approximately
\[
  o_{\text{sat}}
  =0.1v_{\text{the}}+0.7v_{\text{cat}}+0.2v_{\text{sat}}.
\]
Consequently, the updated representation of \texttt{sat} receives most of its
contextual information from \texttt{cat}.

\subsection{Split into attention heads}

In the above example, \(Q\), \(K\), and \(V\) searched for \texttt{cat} using a
verb--subject relationship.

However, in reality multiple attention heads work on multiple such relationships in parallel.

Imagine the sequence:
\begin{quote}
  \texttt{The cat drank milk because it was thirsty.}
\end{quote}
Different relationships may matter at the same time. For example, separate
attention heads might learn to focus on:
\begin{itemize}
  \setlength{\itemsep}{0.25em}
  \item a \textbf{pronoun--noun reference}: \texttt{it} \(\longrightarrow\) \texttt{cat},
  \item a \textbf{verb--subject relationship}: \texttt{drank} \(\longrightarrow\)
        \texttt{cat},
  \item a \textbf{nearby grammatical relationship}: \texttt{the}
        \(\longrightarrow\) \texttt{cat}, and
  \item \textbf{long-range context}: \texttt{thirsty} \(\longrightarrow\) \texttt{cat}.
\end{itemize}

With only one head, a single attention mechanism must represent all these
relationships in the same space. With multiple heads, different heads can
represent different kinds of relationships in parallel:
\begin{center}
\begin{tabular}{@{}ll@{}}
Head 1 & relationship type A \\
Head 2 & relationship type B \\
Head 3 & relationship type C \\
Head 4 & relationship type D
\end{tabular}
\end{center}

The model is not explicitly told what each head should learn. During training,
it discovers useful specializations automatically.

Let
\[
  H=\text{number of heads},
  \qquad
  D=\frac{C}{H}.
\]

For example, if \(C=32\) and \(H=4\), then \(D=8\). \\

So instead of one 32-dimensional attention head, we create four heads, each working with 8 features.


\textbf{Conceptually,} a 32-dimensional query vector
\[
\left[
\begin{array}{cccccccc}
q_1 & q_2 & q_3 & q_4 & q_5 & q_6 & q_7 & q_8 \\
q_9 & q_{10} & q_{11} & q_{12} & q_{13} & q_{14} & q_{15} & q_{16} \\
q_{17} & q_{18} & q_{19} & q_{20} & q_{21} & q_{22} & q_{23} & q_{24} \\
q_{25} & q_{26} & q_{27} & q_{28} & q_{29} & q_{30} & q_{31} & q_{32}
\end{array}
\right]
\]
can be viewed as four smaller query vectors:
\begin{align*}
  \text{Head 1:}&\quad [q_1,\ldots,q_8],\\
  \text{Head 2:}&\quad [q_9,\ldots,q_{16}],\\
  \text{Head 3:}&\quad [q_{17},\ldots,q_{24}],\\
  \text{Head 4:}&\quad [q_{25},\ldots,q_{32}].
\end{align*}
Thus,
\[
  \underbrace{32}_{\text{features}}
  \quad\longrightarrow\quad
  \underbrace{4}_{\text{heads}}
  \times
  \underbrace{8}_{\text{features per head}}.
\]

The shape changes as
\[
  [B,T,C] \longrightarrow [B,H,T,D],
\]
for example,
\[
  [32,8,32] \longrightarrow [32,4,8,8].
\]


\textbf{Meaning:}
\begin{itemize}
  \setlength{\itemsep}{0pt}
  \item 32 sequences in the batch,
  \item 4 attention heads,
  \item 8 token positions, and
  \item 8 features per head.
\end{itemize}
Each head gets its own 8-dimensional slice.

\subsection{Attention scores \( [B,H,T,T] \)}

\textbf{What it does:} Measures how relevant each token is to every other
token.

\[
  S=QK^{\mathsf T},
  \qquad
  S_{\text{scaled}}=\frac{QK^{\mathsf T}}{\sqrt{D}}.
\]

where:
\begin{center}
\begin{tabular}{@{}cl@{}}
  \(Q\) & queries \\
  \(K\) & keys \\
  \(D\) & head size \\
  \(S\) & attention scores
\end{tabular}
\end{center}

For each token's query, we compare it with every token's key using a dot product.

For one attention head in one batch element, the shapes initially are
\[
  Q:[T,D],
  \qquad
  K:[T,D].
\]
Transpose \(K\):
\[
  K^{\mathsf T}:[D,T].
\]
Then matrix multiplication gives
\[
  \underbrace{[T,D]}_{Q}
  \;@\;
  \underbrace{[D,T]}_{K^{\mathsf T}}
  \quad\longrightarrow\quad
  \underbrace{[T,T]}_{\text{attention scores}}.
\]

The dot product measures how aligned two vectors are.

In a real model, the shapes are:
\[
  Q:[B,H,T,D],\qquad
  K^{\mathsf T}:[B,H,D,T]
  \longrightarrow S:[B,H,T,T].
\]

In \(S:[B,H,T,T]\), the dimensions represent:
\begin{center}
\begin{tabular}{@{}cl@{}}
  \(B\) & different sequences in the batch \\
  \(H\) & different attention heads \\
  first \(T\) & query-token positions \\
  second \(T\) & key-token positions
\end{tabular}
\end{center}

\textbf{Example:}

\begin{verbatim}
the cat drank milk because it
\end{verbatim}

When processing \texttt{it}, one attention head might create a query
representing something like:
\begin{quote}
  \emph{``I am looking for the noun this pronoun may refer to.''}
\end{quote}

The keys might produce scores like:
\begin{verbatim}
the      -> 0.2
cat      -> 2.4
drank    -> 0.1
milk     -> 0.5
because  -> 0.0
it       -> 0.7
\end{verbatim}

The model therefore sees:
\[
  \texttt{cat}\longrightarrow\text{strongest match}.
\]

Higher scores indicate stronger relevance.

\begin{quote}
  ``Which tokens matter to me?''
\end{quote}

\subsection{Scaling}

\textbf{What it does:} Stops dot products from becoming too large.

As \(D\) gets larger, dot products tend to grow in magnitude. Large scores can
make softmax too extreme.

\[
  \text{scaled score}=\frac{\text{score}}{\sqrt{D}}.
\]

If \(D=16\), then \(\sqrt{D}=4\), so a score of 12 becomes 3. This helps
softmax behave more smoothly.

\subsection{Causal masking}

\textbf{What it does:} Prevents GPT from seeing future tokens.

\[
  M_{ij}=
  \begin{cases}
    0,       & j\leq i,\\
    -\infty, & j>i,
  \end{cases}
  \qquad
  S'=S+M.
\]

Here,
\begin{center}
\begin{tabular}{@{}cl@{}}
  \(i\) & current query-token position \\
  \(j\) & key-token position being examined
\end{tabular}
\end{center}

If \(j\leq i\), the key token is at the current or a previous position, so attention is allowed. \\
If \(j>i\), the key token is in the future, so attention
is blocked.

For four tokens, with query positions as rows and key positions as columns,
the mask is
\[
  M=
  \begin{bmatrix}
    0 & -\infty & -\infty & -\infty \\
    0 & 0       & -\infty & -\infty \\
    0 & 0       & 0       & -\infty \\
    0 & 0       & 0       & 0
  \end{bmatrix}.
\]

\textbf{Example:}

Suppose the unmasked attention-score matrix is:
\begin{center}
\begin{tabular}{@{}c cccc@{}}
  & \multicolumn{4}{c}{key token} \\
  & \texttt{h} & \texttt{e} & \texttt{l} & \texttt{o} \\
  query \texttt{h} & 1.2 & 0.4 & 0.9 & 0.7 \\
  query \texttt{e} & 1.2 & 0.8 & 1.5 & 0.3 \\
  query \texttt{l} & 0.5 & 1.6 & 1.1 & 0.9 \\
  query \texttt{o} & 0.4 & 0.7 & 1.8 & 1.2
\end{tabular}
\end{center}

Now focus on the row for \texttt{e}:
\[
  [1.2,0.8,1.5,0.3].
\]
Without masking, \texttt{e} could attend to:
\begin{center}
\begin{tabular}{@{}ll@{}}
  \texttt{h} & \\
  \texttt{e} & \\
  \texttt{l} & \(\leftarrow\) future token \\
  \texttt{o} & \(\leftarrow\) future token
\end{tabular}
\end{center}

That is a problem. When GPT is at \texttt{e}, it should not already know that
\texttt{l} and \texttt{o} come later. Therefore, we apply a causal mask. The
masked row becomes
\[
  [1.2,0.8,-\infty,-\infty].
\]

In code, we can first create a lower-triangular mask:
\begin{lstlisting}[style=shadedpython]
mask = torch.tril(
    torch.ones(T, T)
)
\end{lstlisting}
For four tokens, this gives
\[
  \begin{bmatrix}
    1 & 0 & 0 & 0 \\
    1 & 1 & 0 & 0 \\
    1 & 1 & 1 & 0 \\
    1 & 1 & 1 & 1
  \end{bmatrix}.
\]

We then replace every forbidden position with \(-\infty\): \\

\begin{lstlisting}[style=shadedpython]
scores = scores.masked_fill(
    mask == 0,
    float("-inf"),
)
\end{lstlisting}

\subsubsection*{Why use \(-\infty\)?}

The next step is softmax:
\[
  \operatorname{softmax}(s)_i
  =\frac{e^{s_i}}{\sum_j e^{s_j}}.
\]
Since
\[
  e^{-\infty}=0,
\]
masked positions receive zero probability. For example,
\[
  [1.2,0.8,-\infty,-\infty]
  \xrightarrow{\operatorname{softmax}}
  [0.60,0.40,0.00,0.00].
\]
Thus, the query token \texttt{e} gives approximately:
\begin{center}
\begin{tabular}{@{}cl@{}}
  60\% & attention to \texttt{h} \\
  40\% & attention to \texttt{e} \\
  0\%  & attention to \texttt{l} \\
  0\%  & attention to \texttt{o}
\end{tabular}
\end{center}
The future tokens therefore have exactly zero attention probability.

\subsection{Softmax / attention weights}

\textbf{What it does:} Converts scores into probabilities.

\[
  A_{ij}=\frac{e^{S'_{ij}}}{\sum_k e^{S'_{ik}}}.
\]

For example,
\[
  [1.0,2.0,-\infty]
  \xrightarrow{\operatorname{softmax}}
  [0.27,0.73,0.00].
\]
The second token receives approximately 73\% of the attention.

\subsection{Dropout}

\textbf{What it does:} Randomly removes some values during training so the
model does not rely too heavily on particular neurons or attention
connections.

If the dropout probability is \(p\), PyTorch uses \emph{inverted dropout}:
\[
  \operatorname{Dropout}(x)=\frac{m\odot x}{1-p},
\]
where each mask value is sampled independently as
\[
  m_i\sim\operatorname{Bernoulli}(1-p).
\]
Here,
\begin{center}
\begin{tabular}{@{}cl@{}}
  \(x\) & input values \\
  \(p\) & dropout probability \\
  \(m\) & random binary mask containing 0 or 1 \\
  \(\odot\) & element-wise multiplication \\
  \(1-p\) & probability that a value is kept
\end{tabular}
\end{center}

For example, suppose
\[
  p=0.1,
  \qquad
  1-p=0.9,
  \qquad
  x=[0.2,0.6,0.2].
\]
A randomly generated mask might be
\[
  m=[1,0,1].
\]
First apply the mask:
\begin{align*}
  m\odot x
  &=[1,0,1]\odot[0.2,0.6,0.2]\\
  &=[0.2,0,0.2].
\end{align*}
Then divide the surviving values by \(1-p=0.9\):
\[
  \operatorname{Dropout}(x)
  =\frac{[0.2,0,0.2]}{0.9}
  \approx[0.222,0,0.222].
\]
Thus,
\begin{center}
\begin{tabular}{@{}ll@{}}
  Before dropout: & \([0.2,0.6,0.2]\) \\
  Random mask:    & \([1,0,1]\) \\
  After dropout:  & \([0.22,0,0.22]\)
\end{tabular}
\end{center}

Dividing by \(1-p\) keeps the expected magnitude of the values roughly
unchanged during training. With 10\% dropout, each surviving value is scaled
by
\[
  \frac{1}{0.9}\approx1.111,
  \qquad
  0.2\times1.111\approx0.222.
\]

In attention, dropout is often applied after softmax:
\begin{lstlisting}[style=shadedpython]
attention = F.softmax(scores, dim=-1)
attention = self.dropout(attention)
\end{lstlisting}

\subsection{Weighted values \( [B,H,T,D] \)}

\textbf{What it does:} This is the step where attention stops asking
``Who matters?'' and actually collects information from those tokens.

The formula is
\[
  O=AV,
\]
where:
\begin{center}
\begin{tabular}{@{}cl@{}}
  \(A\) & attention weights \\
  \(V\) & value vectors \\
  \(O\) & attention output
\end{tabular}
\end{center}



The full attention formula is
\[
  \boxed{
    \operatorname{Attention}(Q,K,V)
    =\operatorname{softmax}\!\left(
      \frac{QK^{\mathsf T}}{\sqrt{D}}+M
    \right)V
  }.
\]

For attention weights 20\%, 70\%, and 10\%, the output is
\[
  0.2V_A+0.7V_B+0.1V_C.
\]
This is the context-aware representation.

\subsubsection*{Why is it context-aware?}

Before attention, \texttt{sat} had only its own representation. After
attention, its new representation contains information mixed from:
\begin{center}
\begin{tabular}{@{}cl@{}}
  20\% & from \texttt{the} \\
  70\% & from \texttt{cat} \\
  10\% & from \texttt{sat}
\end{tabular}
\end{center}
Thus, \texttt{sat} now carries mostly information from \texttt{cat}.

Therefore, the batched matrix multiplication has shape
\[
  \underbrace{[B,H,T,T]}_{\text{Attention}}
  \;@\;
  \underbrace{[B,H,T,D]}_{\text{Values}}
  \quad\longrightarrow\quad
  \underbrace{[B,H,T,D]}_{\text{Output}}.
\]

\subsection{Combine heads \( [B,T,C] \)}

After each attention head finishes, the output has shape
\[
  [B,H,T,D],
\]
where:
\begin{center}
\begin{tabular}{@{}cl@{}}
  \(B\) & batch size \\
  \(H\) & number of attention heads \\
  \(T\) & sequence length \\
  \(D\) & features per head
\end{tabular}
\end{center}

Since \(H\times D=C\), the heads are concatenated:
\[
  [B,H,T,D] \longrightarrow [B,T,C].
\]

For example, four heads with eight features each give 32 features:
\[
  [B,4,T,8] \longrightarrow [B,T,32].
\]

\subsubsection*{How to combine heads}

First transpose the head and token-position dimensions:
\[
  [B,H,T,D]\longrightarrow[B,T,H,D].
\]
Then flatten the head and per-head feature dimensions:
\[
  [B,T,H,D]\longrightarrow[B,T,H\times D]=[B,T,C].
\]
In code:
\begin{lstlisting}[style=shadedpython]
output = output.transpose(1, 2)
output = output.reshape(B, T, C)
\end{lstlisting}

\subsection{Output projection}

\textbf{What it does:} Concatenation collects the heads. Output projection blends the information from heads.

\subsubsection*{Why do we need this if the heads are already combined?}

Concatenation only places the head outputs next to each other. For example,
suppose
\begin{align*}
  \text{Head 1 output:}&\quad [1,2],\\
  \text{Head 2 output:}&\quad [5,7].
\end{align*}
After concatenation,
\[
  [1,2]\mathbin{\Vert}[5,7]=[1,2,5,7],
\]
where \(\Vert\) denotes concatenation. The information is now together, but
it has not yet been mixed. 

To combine these heads, You apply a learned linear layer:

\[
  Y=OW_O+b.
\]

\begin{lstlisting}[style=shadedpython]
Output projection

[B, T, C]
    |
    v
Linear(C, C)
    |
    v
[B, T, C]
\end{lstlisting}

\begin{quote}
  ``Combine what all the heads learned.''
\end{quote}

\subsection{Residual connection}

\textbf{What it does:} A residual connection means: do not replace the
token's old representation with the new attention result. Instead, add the
attention result to what the token already had.

The formula is:
\[
  x_{\text{new}}=x+\operatorname{Attention}(x).
\]

\textbf{Example:} Suppose a token currently has
\[
  x=[1,2,3],
\]
and attention produces
\[
  \operatorname{Attention}(x)=[0.2,-0.5,0.7].
\]
Then
\begin{align*}
  x_{\text{new}}
  &=[1,2,3]+[0.2,-0.5,0.7]\\
  &=[1.2,1.5,3.7].
\end{align*}
Thus, the old information is still present, but it has been updated.

\begin{quote}
  ``Keep what I already knew and add the update.''
\end{quote}

\subsubsection*{A language example}

Suppose the token is \texttt{sat}. Before attention, its vector might already
encode:
\begin{itemize}
  \setlength{\itemsep}{0pt}
  \item verb-like information,
  \item token identity, and
  \item position information.
\end{itemize}

Attention then looks at the earlier tokens, and discovers that \texttt{cat} is highly relevant, so attention produces an
update containing information about \texttt{cat}. Conceptually, residual
addition performs
\[
  \underbrace{\text{original \texttt{sat} information}}_{x}
  +
  \underbrace{\text{context gathered from \texttt{cat}}}_{\operatorname{Attention}(x)}
  =
  \underbrace{\text{\texttt{sat} understood in this sentence}}_{x_{\text{new}}}.
\]
The result becomes more context-aware without discarding the original token representation. \\

The shapes also remain the same:
\begin{center}
\begin{tabular}{@{}lc@{}}
  \(x\) & \([B,T,C]\) \\
  \(\operatorname{Attention}(x)\) & \([B,T,C]\) \\
  \hline
  \(x_{\text{new}}\) & \([B,T,C]\)
\end{tabular}
\end{center}

\subsection{Feed-forward network}

\begin{lstlisting}[
  style=shadedpython,
  language={},
  basicstyle=\small\sffamily,
  keywordstyle=\bfseries\sffamily,
  morekeywords={ATTENTION,FFN,RESIDUAL,Flow,Details},
  columns=fullflexible,
  keepspaces=true
]
ATTENTION
"Look around."

Which other tokens matter?
Gather their information.

          |
          v

FFN
"Think about what you found."

Take this token's gathered information
and transform or refine it.

Flow:    C -> 4C -> ReLU -> C
Details: input -> more temp features -> switch -ve off -> combine back

          |
          v

RESIDUAL
"Keep what I knew too."

Old representation
+
FFN update
\end{lstlisting}


\subsubsection*{Let us walk through one token}

Suppose that, after attention, the token \texttt{sat} has the representation
\[
  x=[2,1].
\]

\textbf{Step 1---Expand.} The first linear layer transforms the two features
into a larger temporary representation. For example,
\[
  [2,1]\longrightarrow[3,-2,5,-1].
\]
This gives the model more temporary features with which to work.

\textbf{Step 2---Apply ReLU.} ReLU is defined elementwise as
\[
  \operatorname{ReLU}(z)=\max(0,z).
\]
Therefore,
\[
  [3,-2,5,-1]\longrightarrow[3,0,5,0].
\]
The negative values are switched off.

\textbf{Step 3---Shrink back.} The second linear layer combines the four
temporary features back into two:
\[
  [3,0,5,0]\longrightarrow[1.4,2.8].
\]

The entire FFN can therefore be pictured as
\begin{lstlisting}[style=shadedpython,language={}]
Original token representation
[2, 1]
    |
    v  Linear
[3, -2, 5, -1]
    |
    v  ReLU
[3, 0, 5, 0]
    |
    v  Linear
[1.4, 2.8]
\end{lstlisting}

The final vector \([1.4,2.8]\) is the FFN update.

\subsubsection*{Then apply the residual connection}

The old representation is not discarded. The FFN update is added back to it:
\[
  \underbrace{[2,1]}_{\text{old }x}
  +
  \underbrace{[1.4,2.8]}_{\text{FFN update}}
  =
  \underbrace{[3.4,3.8]}_{\text{new }x}.
\]

The sequence

\textbf{Your FFN:}
\[
  x
  \longrightarrow \operatorname{Linear}(C,4C)
  \longrightarrow \operatorname{ReLU}
  \longrightarrow \operatorname{Linear}(4C,C).
\]

\textbf{For this model:}
\[
  32\longrightarrow128
  \longrightarrow\operatorname{ReLU}
  \longrightarrow32.
\]

contains a non-linear ReLU activation. This non-linearity allows the model to learn more complicated patterns than linear transformations alone.

\textbf{FFN Model:}

\[
  x\longrightarrow h\longrightarrow y.
\]

\[
  h=\operatorname{ReLU}(xW_1+b_1),
  \qquad
  y=hW_2+b_2.
\]

\begin{center}
\begin{tabular}{@{}cl@{}}
  \(x\) & what came \emph{into} the FFN \\
  \(h\) & temporary hidden representation (the ``thinking workspace'') \\
  \(y\) & what came \emph{out of} the FFN
\end{tabular}
\end{center}







\(W_1\) : Matrix of the learned weights of the first linear layer,
\(b_1\) : its learned bias. 

In this model,
\[
  W_1:32\longrightarrow128,
  \qquad
  W_1\in\mathbb{R}^{32\times128}.
\]
Thus,
\[
  xW_1+b_1
\]
performs the transformation
\[
  32\text{ numbers}\longrightarrow128\text{ numbers}.
\]

ReLU is then applied to every number:
\begin{center}
\begin{tabular}{@{}ll@{}}
  negative value & \(\longrightarrow 0\) \\
  positive value & \(\longrightarrow\) keep the value
\end{tabular}
\end{center}

The matrix \(W_2\) is another learned weight matrix. It projects the hidden
representation back from \(4C\) features to \(C\) features:
\[
  W_2:4C\longrightarrow C.
\]
For this model, \(W_2\in\mathbb{R}^{128\times32}\), while \(b_2\) is the
bias vector of the second linear layer. The result
\[
  y=hW_2+b_2
\]
is the final FFN output.

\subsubsection*{What are \(W_1\), \(b_1\), \(W_2\), and \(b_2\)?}

They are the trainable parameters of the FFN. In PyTorch,
\begin{lstlisting}[style=shadedpython]
nn.Linear(32, 128)  # W1: weights, b1: bias
nn.Linear(128, 32)  # W2: weights, b2: bias
\end{lstlisting}

During training, gradients flow back to all four parameters:
\begin{lstlisting}[style=shadedpython,language={}]
loss
  |
  v
backpropagation
  |
  v
AdamW
  |
  v
updates W1, b1, W2, and b2
\end{lstlisting}
The model therefore learns how these layers should transform each token's
information.



\subsection{Transformer block}


\begingroup
\small
\renewcommand{\arraystretch}{1.25}
\begin{longtable}{@{}p{0.28\textwidth}p{0.66\textwidth}@{}}
  \textbf{Flow} & \textbf{Description} \\
  \hline
  \endfirsthead
  \textbf{Flow} & \textbf{Description} \\
  \hline
  \endhead

  {\centering\sffamily\bfseries
    \(x\)\par
    \([B,T,C]\)\par}
  &
  \textbf{Step 1: Start with \(x\).}
  At the beginning, each token's knowledge mostly comes from its token
  identity and position:
  \[
    \text{token identity}+\text{position information}.
  \]
  \\

  {\centering\sffamily
    \(\downarrow\)\par
    \textbf{LayerNorm}\par}
  &
  \textbf{Step 2: LayerNorm before attention.}
  Compute
  \[
    \operatorname{LN}(x).
  \]
  LayerNorm stabilizes each token's features: ``Put the values into a stable
  range before attention works on them.''
  \\

  {\centering\sffamily
    \(\downarrow\)\par
    \textbf{Attention}\par
    ``Look around''\par}
  &
  \textbf{Step 3: Attention.}
  Compute
  \[
    \operatorname{Attention}\!\left(\operatorname{LN}(x)\right).
  \]
  Attention lets each token examine the allowed previous tokens and itself.
  For example, \texttt{sat} might assign
  \[
    \texttt{the}:10\%,\qquad
    \texttt{cat}:70\%,\qquad
    \texttt{sat}:20\%.
  \]
  Its attention output is a weighted mixture of those value vectors.
  \\

  {\centering\sffamily
    \(\downarrow\)\par
    \textbf{Attention update}\par
    \([B,T,C]\)\par}
  &
  The attention update has the same shape as \(x\), so the two tensors can be
  added element by element.
  \\

  {\centering\sffamily
    attention update\par
    \(+\) old \(x\)\par
    \(\downarrow\)\par
    \textbf{Residual add}\par
    \([B,T,C]\)\par}
  &
  \textbf{Step 4: First residual connection.}
  The update is
  \[
    x\leftarrow x+
    \operatorname{Attention}\!\left(\operatorname{LN}(x)\right).
  \]
  In words,
  \[
    \text{old }x+\text{attention update}=\text{new }x.
  \]
  \\

  {\centering\sffamily
    \(\downarrow\)\par
    \textbf{LayerNorm}\par}
  &
  \textbf{Step 5: LayerNorm again.}
  The updated \(x\) passes through
  \[
    \operatorname{LN}(x),
  \]
  which prepares the representation for the FFN.
  \\

  {\centering\sffamily
    \(\downarrow\)\par
    \textbf{FFN}\par
    ``Think about what you found''\par
    \(C\rightarrow4C\rightarrow\operatorname{ReLU}\rightarrow C\)\par}
  &
  \textbf{Step 6: FFN.}
  Compute
  \[
    \operatorname{FFN}\!\left(\operatorname{LN}(x)\right).
  \]
  The FFN works on each token independently. If \(C=32\), it typically
  performs
  \[
    32\longrightarrow128
    \longrightarrow\operatorname{ReLU}
    \longrightarrow32.
  \]
  For \texttt{sat},
  \[
    \text{context-aware \texttt{sat} vector}
    \longrightarrow\operatorname{FFN}
    \longrightarrow\text{refined \texttt{sat} vector}.
  \]
  Attention has already gathered information from \texttt{cat}; the FFN now
  transforms the information stored inside \texttt{sat}.
  \\

  {\centering\sffamily
    \(\downarrow\)\par
    \textbf{FFN update}\par
    \([B,T,C]\)\par}
  &
  The FFN projects back to \(C\) features, so its update again has shape
  \([B,T,C]\).
  \\

  {\centering\sffamily
    FFN update\par
    \(+\) old \(x\)\par
    \(\downarrow\)\par
    \textbf{Residual add}\par
    \([B,T,C]\)\par}
  &
  \textbf{Step 7: Second residual connection.}
  The second update is
  \[
    x\leftarrow x+
    \operatorname{FFN}\!\left(\operatorname{LN}(x)\right).
  \]
  Again,
  \[
    \text{old }x+\text{FFN update}=\text{new }x.
  \]
  This keeps the contextual representation and adds the FFN's refinement.
  \\

  {\centering\sffamily
    \(\downarrow\)\par
    \textbf{Block output}\par
    \([B,T,C]\)\par}
  &
  The block produces a more context-aware and refined representation while
  preserving the original tensor shape.
  \\
\end{longtable}
\endgroup



\subsection{Repeat Transformer blocks}

Now we take one Transformer block and stack many of them.

If there are \(N\) layers,
\[
  X_N=\operatorname{Block}_N\!\left(
    \cdots\operatorname{Block}_2\!\left(
      \operatorname{Block}_1(X)
    \right)
  \right).
\]

Each layer refines the representation, while its shape remains \([B,T,C]\).

Each block receives
\[
  [B,T,C]
\]
and returns
\[
  [B,T,C].
\]
The shape stays the same, but the meaning encoded by the numbers becomes
richer.



\begin{lstlisting}[
  style=shadedpython,
  language={},
  basicstyle=\small\sffamily,
  keywordstyle=\bfseries\sffamily,
  morekeywords={Initial,Block},
  columns=fullflexible,
  keepspaces=true
]
Initial representation
        |
        v
rough understanding

Block 1
        |
        v
better understanding

Block 2
        |
        v
more contextual understanding

Block 3
        |
        v
even richer understanding
\end{lstlisting}

In code, the stack can look like this:
\begin{lstlisting}[style=shadedpython]
self.blocks = nn.Sequential(
    *[
        TransformerBlock()
        for _ in range(number_of_layers)
    ]
)
\end{lstlisting}
Then the representation is passed through every block in sequence:
\begin{lstlisting}[style=shadedpython]
x = self.blocks(x)
\end{lstlisting}



\subsection{Final LayerNorm}

\[
  X_{\text{final}}=\operatorname{LN}(X_N),
  \qquad
  X_{\text{final}}:[B,T,C].
\]

It prepares the final features for prediction.

\subsection{Language-model head \( [B,T,V] \)}

\textbf{What it does:} The Language Model Head is the final layer that converts the Transformer's internal token representation into:

one score for every possible token in the vocabulary.

\[
  Z=XW_{\text{vocab}}+b,
  \qquad
  [B,T,C]\longrightarrow[B,T,V].
\]

\begin{center}
\begin{tabular}{@{}cl@{}}
  \(X\) & current token representations \\
  \(W_{\text{vocab}}\) & learned matrix mapping \(C\) features to \(V\) vocabulary scores \\
  \(b\) & learned bias \\
  \(Z\) & final logits
\end{tabular}
\end{center}



This means:
\begin{lstlisting}[
  style=shadedpython,
  language={},
  basicstyle=\small\sffamily,
  keywordstyle=\bfseries\sffamily,
  morekeywords={Transformer,Linear,Vocabulary},
  columns=fullflexible,
  keepspaces=true
]
Transformer representation
[B, T, C]
       |
       v
Linear layer
C -> V
       |
       v
Vocabulary scores
[B, T, V]
\end{lstlisting}



For one token, imagine that the Transformer produced
\[
  x=[0.2,-0.7,1.1,\ldots,0.4],
  \qquad x\in\mathbb{R}^{32}.
\]
These 32 numbers are useful internal features, but they are not token
probabilities yet.

Suppose the vocabulary has \(V=10\) tokens:
\begin{center}
\begin{tabular}{@{}cl@{}}
  0 & \texttt{space} \\
  1 & \texttt{d} \\
  2 & \texttt{e} \\
  3 & \texttt{h} \\
  4 & \texttt{l} \\
  5 & \texttt{o} \\
  6 & \texttt{p} \\
  7 & \texttt{r} \\
  8 & \texttt{t} \\
  9 & \texttt{w}
\end{tabular}
\end{center}
The language-model head converts the 32 features into 10 vocabulary scores:
\begin{lstlisting}[
  style=shadedpython,
  language={},
  basicstyle=\small\sffamily,
  keywordstyle=\bfseries\sffamily,
  morekeywords={Transformer,Language,Model,Head},
  columns=fullflexible,
  keepspaces=true
]
Transformer output for one position
[0.2, -0.7, 1.1, ..., 0.4]
            32 features
                 |
                 v
        Language Model Head
                 |
                 v
space -> -0.5
d     -> -1.2
e     ->  0.8
h     -> -0.9
l     ->  1.2
o     ->  3.1
p     -> -0.4
r     ->  0.6
t     ->  0.1
w     ->  1.0
\end{lstlisting}

These scores are called \emph{logits}. The highest logit here is
\[
  \texttt{o}\longrightarrow3.1,
\]
so the model currently considers \texttt{o} the strongest candidate for the
next token.

During generation, softmax converts all vocabulary logits into
probabilities. \\

An illustrative partial distribution might be:

\begin{lstlisting}[style=shadedpython,language={}]
logits                    probabilities
space -> -0.5             space ->  2%
e     ->  0.8             e     ->  8%
l     ->  1.2             l     -> 12%
o     ->  3.1             o     -> 55%
...                        ...
\end{lstlisting}
The probabilities across the complete vocabulary sum to 100\%. \\

The model can then sample the next token from this distribution. \\

\begin{quote}
  ``What token should come next?''
\end{quote} 


\subsection{Cross-entropy loss}

These three steps form one continuous training loop:
\begin{lstlisting}[
  style=shadedpython,
  language={},
  basicstyle=\small\sffamily,
  keywordstyle=\bfseries\sffamily,
  morekeywords={Prediction,Cross,Backpropagation,AdamW},
  columns=fullflexible,
  keepspaces=true
]
Prediction
     |
     v
Cross-entropy loss
"How wrong was I?"
     |
     v
Backpropagation
"Which weights caused that error,
 and in what direction?"
     |
     v
AdamW
"Update those weights to reduce
 the error next time."
\end{lstlisting}

After the language-model head, the tensor shapes are
\begin{lstlisting}[style=shadedpython]
logits.shape  = [B, T, V]
targets.shape = [B, T]
\end{lstlisting}

Suppose
\[
  B=32,\qquad T=8,\qquad V=10.
\]
Then
\begin{lstlisting}[style=shadedpython]
logits.shape  = [32, 8, 10]
targets.shape = [32, 8]
\end{lstlisting}

For every token position, the model produces 10 scores because there are 10
possible vocabulary tokens. At one position, those scores correspond to
\begin{lstlisting}[style=shadedpython,language={}]
Possible next tokens:

space -> score
d     -> score
e     -> score
h     -> score
l     -> score
o     -> score
p     -> score
r     -> score
t     -> score
w     -> score
\end{lstlisting}
The target for that position contains only one value:
the ID of the single correct next token.

\subsubsection*{Why flatten?}

There are
\[
  \underbrace{32}_{\text{sequences}}
  \times
  \underbrace{8}_{\text{positions per sequence}}
  =256
\]
separate next-token predictions.

Cross-entropy expects one row of vocabulary scores for each prediction.
Because \(V=10\), this means 256 rows, each containing 10 scores:
\[
  [B,T,V]\longrightarrow[B\times T,V].
\]
Numerically,
\[
  [32,8,10]\longrightarrow[256,10].
\]

The targets are flattened in the same way:
\[
  [B,T]\longrightarrow[B\times T],
  \qquad
  [32,8]\longrightarrow[256].
\]

Thus, after flattening,
\begin{lstlisting}[style=shadedpython]
logits.shape  = [256, 10]
targets.shape = [256]
\end{lstlisting}
There are 256 prediction problems, and each problem chooses among 10
vocabulary tokens.

For one prediction,
\[
  L=-\log p(y),
\]
where \(y\) is the correct token. 

If the correct token receives 90\% probability,
\[
  -\log(0.9)\approx0.105.
\]
If it receives 1\% probability,
\[
  -\log(0.01)\approx4.605.
\]

\begin{quote}
  ``How wrong was my prediction?''
\end{quote}

\subsubsection*{Many predictions at once}

In our example, there are 256 predictions. Suppose their individual losses
were
\begin{lstlisting}[style=shadedpython,language={}]
prediction 1   -> 0.2
prediction 2   -> 0.6
prediction 3   -> 1.1
...
prediction 256 -> 0.3
\end{lstlisting}

PyTorch usually averages these individual losses:
\[
  L=-\frac{1}{N}\sum_{i=1}^{N}\log p(y_i),
\]
where
\[
  N=B\times T.
\]
Here,
\[
  N=32\times8=256.
\]

The result is one final scalar, for example,
\[
  \text{loss}=0.82.
\]
This single number summarizes the model's average prediction error across all
256 token positions.

\subsection{Backpropagation}

Suppose the final loss is
\[
  \text{loss}=0.82.
\]
We know that the model made errors, but it may contain thousands or millions
of parameters, including
\begin{itemize}
  \setlength{\itemsep}{0pt}
  \setlength{\parsep}{0pt}
  \item token-embedding weights,
  \item position-embedding weights,
  \item query, key, and value weights,
  \item FFN weights, and
  \item language-model-head weights.
\end{itemize}

Which parameters should change, and by how much? That is what
backpropagation calculates.

For every parameter \(\theta\), it computes
\[
  \frac{\partial L}{\partial\theta}.
\]
This quantity is called the \emph{gradient}. \\ 

It measures how the loss changes when that parameter changes.

For example, if a weight is 0.5 and its gradient is \(+0.2\), the positive
gradient indicates the direction that increases the loss, so the optimizer
moves in the opposite direction.

\begin{quote}
  ``Which parameters caused the error?''
\end{quote}

\subsubsection*{Backpropagation does not update the weights}

This is an important distinction: calling
\begin{lstlisting}[style=shadedpython]
loss.backward()
\end{lstlisting}
does not change the model parameters. It only calculates their gradients.

Conceptually, \texttt{loss.backward()} means:
\begin{quote}
  ``Calculate how every trainable parameter affected this loss.''
\end{quote}

For example, backpropagation might calculate
\begin{lstlisting}[style=shadedpython,language={}]
weight A gradient = +0.20
weight B gradient = -0.07
weight C gradient = +0.91
weight D gradient = -0.13
...
\end{lstlisting}
These gradients are stored in each parameter's \texttt{.grad} attribute. \\

The optimizer uses them in the next step to update the weights.



\subsection{AdamW update}

A simplified gradient-descent view is
\[
  \theta_{\text{new}}
  =\theta_{\text{old}}
  -\eta\frac{\partial L}{\partial\theta},
\]

The derivative
\[
  \frac{\partial L}{\partial\theta}
\]
describes how the loss changes when \(\theta\) changes.

If the gradient is negative, increasing \(\theta\) decreases the loss. The
update therefore subtracts a negative number:
\[
  \frac{\partial L}{\partial\theta}<0
  \quad\Longrightarrow\quad
  \theta_{\text{new}}
  =\theta_{\text{old}}
  +\eta\left\lvert\frac{\partial L}{\partial\theta}\right\rvert.
\]
Thus, \(\theta\) increases in the direction that reduces the loss. \\ 
If the gradient is positive, the subtraction makes \(\theta\) smaller. 
In both cases, gradient descent moves opposite to the gradient.

where \(\eta\) is the learning rate. \\


\subsubsection*{Why AdamW instead of basic gradient descent?}

The formula above is a simplified explanation. AdamW is more sophisticated
than the basic update
\begin{lstlisting}[style=shadedpython,language={}]
weight = weight - learning_rate * gradient
\end{lstlisting}
because it also tracks recent gradient behavior, adjusts the effective step
size, and applies weight decay.

\textbf{Momentum.} If gradients repeatedly point in the same direction,
\begin{lstlisting}[style=shadedpython,language={}]
step 1 -> right
step 2 -> right
step 3 -> right
\end{lstlisting}
Adam builds momentum in that direction. This can make optimization faster and
less sensitive to noisy changes in individual gradients.

\textbf{Adaptive scaling.} Some parameters may have large gradients, such as
\[
  \text{gradient}=10,
\]
while others may have tiny gradients, such as
\[
  \text{gradient}=0.001.
\]
Adam adjusts the effective step size for each parameter using its history of
gradient magnitudes.

\textbf{Weight decay.} AdamW also gently discourages weights from becoming
unnecessarily large. Its decoupled weight decay is applied separately from
the gradient-based update.

Mathematically, the learning process can be summarized as
\[
  \text{prediction}
  \longrightarrow
  L=-\log p(y)
  \longrightarrow
  \frac{\partial L}{\partial\theta}
  \longrightarrow
  \theta_{\text{new}}
  =\theta_{\text{old}}
  -\eta\frac{\partial L}{\partial\theta}.
\]
This is essentially the heart of how your MiniGPT learns.



\subsection{Generation}

Suppose the current text is \texttt{hello}. Run the tokenized context through
the model:
\[
  \text{tokens }[B,T]
  \longrightarrow
  \text{logits }[B,T,V].
\]

Keep only the final position:
\begin{verbatim}
latest_logits = logits[:, -1, :]
\end{verbatim}
which has shape \([B,V]\). Convert these logits into probabilities:
\[
  p_i=\frac{e^{z_i}}{\sum_j e^{z_j}}.
\]

\textbf{Example:}
\begin{verbatim}
" " -> 60%
"w" -> 20%
"e" -> 10%
\end{verbatim}

If a space is sampled, append it to \texttt{hello}, then repeat:
\begin{verbatim}
hello
-> hello 
-> hello w
-> hello wo
-> hello wor
-> ...
\end{verbatim}

\begin{quote}
  \textbf{Predict \(\rightarrow\) sample \(\rightarrow\) append
  \(\rightarrow\) repeat.}
\end{quote}

\subsection*{One compact mathematical pipeline}

\[
  \text{tokens}
  \longrightarrow
  X=E_{\text{token}}+E_{\text{position}}
\]
\[
  \downarrow
\]
\[
  X\leftarrow X+
  \operatorname{Attention}\!\left(\operatorname{LN}(X)\right)
\]
\[
  \downarrow
\]
\[
  X\leftarrow X+
  \operatorname{FFN}\!\left(\operatorname{LN}(X)\right)
\]
\[
  \downarrow
\]
\[
  Z=XW_{\text{vocab}}+b
  \longrightarrow
  P=\operatorname{softmax}(Z).
\]

During training,
\[
  P\longrightarrow\text{cross-entropy}
  \longrightarrow\text{backpropagation}
  \longrightarrow\text{AdamW}.
\]

During generation,
\[
  P\longrightarrow\text{sample next token}
  \longrightarrow\text{append}
  \longrightarrow\text{repeat}.
\]

The central Transformer formula is
\[
  \boxed{
    \operatorname{Attention}(Q,K,V)
    =\operatorname{softmax}\!\left(
      \frac{QK^{\mathsf T}}{\sqrt{D}}+M
    \right)V
  }.
\]

\section{What you have now built}

This is a real decoder-only Transformer. Its scale and implementation details
differ from GPT-style models, but the core algorithmic idea is the same.

\begin{center}
\renewcommand{\arraystretch}{1.25}
\begin{tabular}{@{}p{0.27\textwidth}p{0.28\textwidth}p{0.33\textwidth}@{}}
  \textbf{Component} & \textbf{Your model} & \textbf{GPT-style model} \\
  \hline
  Tokenizer & Character-level & BPE or another subword tokenizer \\
  Embeddings & Yes & Yes \\
  Position encoding & Learned position embeddings & Usually RoPE in modern models \\
  Attention & Causal multi-head attention & Causal multi-head attention \\
  Transformer blocks & 4 & Dozens or hundreds \\
  Parameters & Roughly \(100\text{K}\)--\(1\text{M}\) & Billions in large models \\
  Training data & A small text dataset & Up to trillions of tokens \\
  Objective & Next-token prediction & Next-token prediction
\end{tabular}
\end{center}

The model sizes are dramatically different, but both learn by predicting the
next token from the tokens that came before it.
